\documentclass[11pt,a4paper]{article}
\usepackage{brainteaser}

\begin{document}
\brainteasertitle{1}{Jeu des entiers avec pénalité}

\begin{statementbox}
Deux joueurs choisissent simultanément un entier $n_1,n_2$ entre $0$ et
$100$. On retire $10$ au plus grand des deux nombres, puis on compare les
valeurs obtenues. La plus grande valeur gagne; une égalité vaut $0$.

Le joueur J2 connaît la stratégie aléatoire de J1, mais pas le nombre
effectivement tiré. Quelle stratégie J1 doit-il adopter pour maximiser sa
garantie de gain?
\end{statementbox}

\begin{hintsbox}
\begin{enumerate}
  \item Poser $d=n_1-n_2$. Pour quelles valeurs de $d$ J1 gagne-t-il, perd-il
  ou fait-il égalité?
  \item Pourquoi aucun choix déterministe de J1 ne peut-il être optimal?
  \item Chercher un bloc de nombres sur lequel chaque choix adverse rencontre
  le même nombre de victoires, de défaites et d'égalités.
  \item Pourquoi un bloc de longueur $20$ apparaît-il naturellement?
  \item Pourquoi faut-il placer ce bloc contre la borne supérieure $100$?
  \item Pour prouver l'optimalité, imaginer que J2 joue uniformément sur ce
  même bloc.
\end{enumerate}
\end{hintsbox}

\correction

J1 gagne lorsque
\[
n_1-n_2\geq 11
\qquad\text{ou}\qquad
1\leq n_2-n_1\leq 9.
\]
Il perd dans les situations symétriques et fait égalité lorsque les nombres
sont égaux ou séparés de $10$.

\begin{keyideabox}
J1 choisit uniformément parmi les vingt nombres
\[
S=\{81,82,\ldots,100\}.
\]
\end{keyideabox}

Contre un choix fixé $n_2\in S$, les vingt choix possibles de J1 se
répartissent toujours en
\[
9\text{ victoires},\qquad 9\text{ défaites},\qquad 2\text{ égalités}.
\]
Ainsi,
\[
\Pp(\text{victoire})=\frac9{20},\qquad
\Pp(\text{défaite})=\frac9{20},\qquad
\Pp(\text{égalité})=\frac1{10}.
\]

Si J2 choisit un nombre inférieur à $81$, la probabilité de victoire de J1
est au moins $1/2$. Le pire cas est donc bien atteint dans $S$.

Pour prouver qu'on ne peut garantir davantage, supposons que J2 joue lui aussi
uniformément sur $S$. Tout nombre fixe de J1 bat au plus neuf des vingt nombres
de J2: exactement neuf s'il appartient à $S$, et au plus neuf s'il est
inférieur à $81$. Une combinaison aléatoire ne peut pas dépasser cette borne.

\begin{resultbox}
La stratégie minimax est la loi uniforme sur $\{81,\ldots,100\}$. Elle garantit
$45\%$ de victoires, $45\%$ de défaites et $10\%$ d'égalités. Avec les gains
$+1,-1,0$, le gain moyen garanti vaut $0$.
\end{resultbox}

\end{document}
